\section{Beyond Debt, Never Beneath Justice}

Friendship with God does not mark a graduation from rule into moral spontaneity. Charity perfects justice rather than superseding it: love becomes mature not by escaping the good but by willing it as the good of the beloved.

Aquinas calls charity the \emph{form of the virtues}: it directs their acts toward the last end and gives them the form of acts belonging to friendship with God while each virtue retains its proper object (ST II--II, q.~23, a.~8). Justice still renders what is due. Religion still renders worship to the first principle. Charity gives those acts their final Godward orientation (CCC 1827, 2095).

Because worship does not repair a divine lack, justice recognizes the Creator's excellence and the creature's received condition, while charity loves the One justice acknowledges. A person can therefore render what is due as an act of love, obey without mercenary calculation, and choose the act freely without making it optional.

Aquinas names peace as charity's proper effect. Justice causes peace indirectly by removing injuries and other obstacles; charity causes it directly by unifying desire in God and joining persons in one good (ST II--II, q.~29, a.~3). The peace of the wayfarer remains imperfect and can coexist with disturbance, dryness, and conflict. Its formal claim is not that the person always feels settled, but that desire need no longer be organized around resistance to the good.

Charity also exceeds the calculation of repayment. A debt can be treated as discharged once an equal measure has been returned. No creature can equal the gift of being, grace, or promised beatitude, and God has no need of compensation. Minimal compliance looks for the point at which the account is settled; love attends instead to the good of the beloved. Determinate obligations remain, but the horizon within which they are performed has changed.

The neighbor is the necessary test. Paul says that love fulfills the law because it does no evil to the neighbor (Rom 13:8--10). The First Letter of John is more severe: one cannot claim love of the unseen God while hating the brother who is seen (1 John 4:20--21). Any mystical language that ends in fascination with one's own interior state has already failed that test.

Augustine's compressed instruction, ``Love, and do what you will,'' is often made to say the opposite of its context. In his seventh homily on First John, outwardly similar acts can arise from contrary roots: a harsh correction may seek a person's good, while a caress may conceal predation. Augustine does not authorize impulse. He demands that silence, speech, correction, and mercy proceed from charity, because an act is morally disclosed by the love from which it comes (\emph{Homilies on First John} VII.8--9).

That root has objective form. Charity wills true good; it does not call evil good in order to appear gentle. It can correct without hatred, forgive without denying injury, and suffer with another without making suffering desirable. The Cross governs the meaning of charity. Divine love is manifested not as indiscriminate affirmation but in the Father sending the Son so that sinners may live, and in the Son freely laying down his life for those who did not first love him (1 John 4:9--11; John 10:17--18).

Charity takes both contemplative and ordinary form. To adore, keep faith, tell truth, restore what is owed, feed the hungry, forgive an enemy, bear another's burden, and accept correction can all be informed and commanded by charity. These finite acts are precisely where the love of God is tested: the received good of another person makes a claim upon my will.

Justice answers: \emph{this good is due}. Charity answers: \emph{I will it for the beloved in God}.

Love goes beyond the calculation of debt. It never falls beneath justice.
